Calcula el \textbf{volumen}, el \textbf{área lateral} y el \textbf{área total} de las siguientes figuras:

			\begin{parts}
				\part Cilindro con altura $h=17$ cm y un radio $r=4$ cm.

				\begin{minipage}{.25\linewidth}
					\includegraphics[width=0.8\linewidth]{mex_0024.png}
				\end{minipage}\hfill
				\begin{minipage}{.75\linewidth}
					Volumen:\\[-2em]
					\begin{solutionbox}{1cm}
						\[V=\pi r^2h=(3.14) 4^2\cdot 17= 857.12\]
					\end{solutionbox}
					
					A. Lateral:\\[-2em]
					\begin{solutionbox}{1cm}
						\[A_L=2\pi rh=2(3.14) 4\cdot 17= 2(3.14) 68=428.48\]
					\end{solutionbox}
					
					A. Total:\\[-2em]
					\begin{solutionbox}{1cm}
						\[A_T=A_L+2\pi r^2=428.48+2(3.14) 16=528.96\]
					\end{solutionbox}
				\end{minipage}

				\part Prisma octagonal de 19 cm de altura y su base es un octágono cuyos lados $l$ miden 7 cm y un apotema $a$ de 5 cm.
				
				\begin{minipage}{.25\linewidth}
					\includegraphics[width=0.8\linewidth]{mex_0025.png}
				\end{minipage}\hfill
				\begin{minipage}{.75\linewidth}
					Volumen:\\[-2em]
					\begin{solutionbox}{1cm}
						\[V=A_b\cdot h= \left(\dfrac{nla}{2}\right)h= \dfrac{8(7)5}{2}(19)=2660\]
					\end{solutionbox}
					
					A. Lateral:\\[-2em]
					\begin{solutionbox}{1cm}
						\[A_L=nlh=8(7)19= 1064\]
					\end{solutionbox}
					
					A. Total:\\[-2em]
					\begin{solutionbox}{1cm}
						\[A_T=A_L+2\dfrac{nla}{2}=A_L+nla=1064+280=1344\]
					\end{solutionbox}
				\end{minipage}
			\end{parts}